\pg{348} b) Os entes finitos, ainda que todos possuam uma realidade ou perfeição limitada, todavia nem todos a possuem no mesmo grau, mas um é mais perfeito, e outro o é menos; assim o ente vivo é mais perfeito do que o ente nàáo-vivo, Tal desegualdade deve ter a sua causa sufficiente. Esta causa deve encontrar-se --- ou nos proprios entes finitos, --- ou num ente infinito. Ora não pode encontrar-se nos entes finitos; porque estes, sendo contingentes, não tém em si a razão da sua existencia, e por isso nem do grau da sua realidade ou perfeição. Logo deve encontrarse num ente infinito, que, espalhando diversamente as perfeições em conformidade com os archetypos da sua intelligencia, produziu a maravilhosa variedade do universo. Logo a existencia dos entes \emph{finitos} suppõe a existencia do Ente \emph{infinito} \footnote{\emph{No} mundo ha muitas e diversas perfeigóes, mas todas limitadas e num grau desegual; assim, no mesmo genero de seres vivos, o homem é mais perfeito do que o animal, e este é mais perfeito do que a planta. Ora tal limitação e desegualdade, que se encontra na perfeição das creaturas, não pode derivar da essencia ou natureza d'ellas. Na verdade, segundo S. Thomaz, uma perfeição.. que convem-a um ente pela sua essencia, não pode ser de modo algum limitada, ou diminuta, sem que a essencia deixe de ser o que é, mas deve existir nelle no grau mais elevado, de que a propria perfeição é susceptivel. Logo essa limitação e desegualdade deve derivar, em ultima analyse, de uma causa extrinseca, que possua por si todas as perfeições num grau eminente'e que as communique por differentes modos aos entes creados (\emph{Gen}£., II, 12). --- O S. Doutor ensina tambem que, quando uma coisa se diz e é mais perfeita que outra, essas coisas dizem-se e são mais ou menos perfeitas pelo facto de dependerem e de se approximarem, mais ou menos, de um ente unico e perfeitissimo: porque, se cada coisa possuisse \emph{por si} mesma a perfeição, de que é dotada, não haveria razão por que uma devesse ser mais e outra menos perfeita (\emph{De pot}., q. 3, a. 9).}.

71. Corollrios. --- Das doutrinas expostas resulta :

I. O infinito presuppoe-se ao finito --- como causa, porque o infinito --produz-o-finito;-----mas-nàao como \emph{objecto}, porque o infinito é por nós conhecido.depois e por meio do finito (\textasciicircum{}).\footnote{A existencia \emph{real} do ente infinito não pode ser demonstrada apenas pelo \emph{conceito} do mesmo ente. \textbar{} certo que nada pode conceber-se que seja mais perfeito do que o ente \emph{infinito}, e que a existencia é uma perfeição ; mas, emquanto não se sahe do conceito, a existencia, devida ao Ente infinito, será apenas \emph{ideal}. É o que dissemos em relação ao ente \emph{necessario}.}

